\chapter{Determinants as Volume}\label{linear:ch:volume}

\begin{center}\itshape
``Then you should say what you mean,'' the March Hare went on.\\
``I do,'' Alice hastily replied; ``at least---at least I mean what I say---that's the same thing, you know.''
\upshape(Lewis Carroll, \emph{Alice's Adventures in Wonderland}, 1865)
\end{center}
\noindent\emph{Use Volume~1's linear algebra and determinant uniqueness theorem. The identities here support Appendices~\ref{linear:ch:eigen} and~\ref{linear:ch:inner}; Volume~9 uses the area and orientation interpretation.}

Queen Victoria supposedly enjoyed \emph{Alice} so much that she requested its author's next book and received a treatise on determinants. Dodgson himself denied the story. The treatise was real, though, as was his 1866 Royal Society paper on condensation: repeatedly replace adjacent $2\times2$ blocks by their determinants, dividing by an entry from the preceding stage. With nonzero divisors this small operation recovers the full determinant exactly.

Robbins and Rumsey revisited the rule in the early 1980s, replacing $ad-bc$ by $ad+\lambda bc$. Their expansion involved alternating sign matrices, whose entries are $0,1,-1$ and whose nonzero entries alternate along each row and column. The counts begin $1,2,7,42,429,7436$. Mills, Robbins and Rumsey conjectured their product formula in 1983; Zeilberger proved it in 1996, and Kuperberg gave a different proof that year through statistical mechanics.

Signed volume provides the algebraic starting point. Exterior powers collect the minors measuring lower-dimensional projections; composition then gives Cauchy--Binet. Complementary minors lead to Jacobi's identity and condensation. Finally, a tail-swapping cancellation turns determinants of path counts into counts of disjoint paths, and hence into lozenge-tiling counts.

\begin{think}
A regular lattice hexagon of side two contains $24$ unit triangles. Pair adjacent triangles into lozenges: there are $20$ tilings. The side-three count is $980$, while side ten already gives a $34$-digit integer. Try drawing a few tilings as stacks of cubes. Section~\ref{linear:sec:lgv} converts these counts into binomial determinants; the small cases can then be evaluated exactly.
\end{think}

\section{Volume and orientation}

Vectors $v_1,\ldots,v_n\in\R^n$ span the parallelepiped $\{\sum_it_iv_i:0\le t_i\le1\}$. Geometric volume suggests three rules: stretching one edge by $\lambda\ge0$ multiplies it by $\lambda$; a shear adding one edge to a multiple of another preserves it; and the coordinate unit cube has volume one. Figure~\ref{linear:fig:shear} shows the planar shear.

To allow negative scale factors, keep track of orientation and use signed volume. We take multilinearity and alternation as the precise algebraic axioms suggested by these rules; the geometric rules alone have not been used as a proof of multilinearity. Volume~1, Theorem~7.3, says the unique alternating multilinear form normalised at the coordinate basis is
\[D(v_1,\ldots,v_n)=\det(v_1,\ldots,v_n).
\]
\begin{figure}[ht]
\centering
\fig{linear-ch02-shear}
\caption{Replacing $v_2$ by $v_2+v_1$ slides the parallelogram along its base: same base, same height, same area. In
coordinates the area is $|\det(v_1,v_2)|$.}
\label{linear:fig:shear}
\end{figure}

\begin{remark}
Volume~7, Chapter~9, proves that a linear map scales Lebesgue volume by $|\det T|$, making the geometric interpretation rigorous. Here determinant algebra is the premise; volume is its interpretation.
\end{remark}

The sign of a nonzero determinant distinguishes the two orientations of a real basis. In the plane the positive orientation turns counterclockwise from the first vector to the second.

\begin{example}[the shoelace formula]\label{linear:ex:shoelace}
For a simple polygon with counterclockwise vertices $p_k=(x_k,y_k)$, $p_{m+1}=p_1$, its area is
\[\operatorname{area}(P)=\frac12\sum_{k=1}^m\det(p_k,p_{k+1})
 =\frac12\sum_{k=1}^m(x_ky_{k+1}-x_{k+1}y_k).
\]
Each summand is the signed area of the triangle from zero to an edge. For a convex polygon containing zero these triangles partition it. In general their signed coverage has multiplicity one inside and zero outside, as given by the boundary's winding number. Integrating that multiplicity gives the same formula, including nonconvex simple polygons. Reversing vertex order reverses its sign; self-intersecting polygons instead give signed winding area. This is why the simplicity and orientation conventions matter in a polygon-area computation.
\end{example}

\section{Exterior powers}\label{linear:sec:exterior}

For $k$ vectors in a higher-dimensional space, projected signed volumes form a whole list of minors. An exterior power stores that list as one vector.

Choose a basis $e_1,\ldots,e_n$ of $V$. For $1\le k\le n$, form a space $\Lambda^kV$ with basis symbols $e_I$, one for each increasing $k$-subset $I$ of $[n]$. If $M$ is the coordinate matrix of $v_1,\ldots,v_k$, define
\[v_1\wedge\cdots\wedge v_k=\sum_{|I|=k}\det(M_I)e_I.
\]
Here $M_I$ takes rows $I$. Also set $\Lambda^0V=\mathbb K$ and $\Lambda^kV=0$ for $k>n$, with empty determinants equal to one. These conventions include the boundary cases of the identities below.

At $k=1$ the construction returns the original vector; at $k=n$ it is the determinant times $e_{[n]}$. For $k=2,n=3$, its coordinates are $v_1w_2-v_2w_1$, $v_1w_3-v_3w_1$ and $v_2w_3-v_3w_2$, the cross-product coordinates up to order and sign. Their Euclidean squared sum will give parallelogram area squared in Exercise~\ref{linear:ex:lagrangeid}.

\begin{theorem}[exterior powers]\label{linear:thm:exterior}
\begin{enumerate}[label=(\alph*)]
\item The exterior product is alternating and multilinear, with $e_{i_1}\wedge\cdots\wedge e_{i_k}=e_I$ for increasing indices.
\item Every alternating multilinear $\varphi:V^k\to U$ factors uniquely through a linear $\tilde\varphi:\Lambda^kV\to U$, so $\varphi(v_1,\ldots,v_k)=\tilde\varphi(v_1\wedge\cdots\wedge v_k)$.
\item A family $v_1,\ldots,v_k$ is independent precisely when its exterior product is nonzero.
\end{enumerate}
\end{theorem}
\begin{proof}
(a) Each minor is alternating and multilinear in the columns. Coordinate basis columns give one identity minor and zero for every other row set.

(b) The basis forces $\tilde\varphi(e_I)=\varphi(e_{i_1},\ldots,e_{i_k})$. Define it that way. Expanding each $v_j=\sum_im_{ij}e_i$, repeated indices vanish by alternation. The injective index choices group by their image $I$ and permutation. Their signed coefficient sum is $\det(M_I)$, so the expansion is exactly $\sum_I\det(M_I)\tilde\varphi(e_I)$. This proves existence as well as uniqueness, including in characteristic two, where alternation still means vanishing on repeated arguments.

(c) Dependence gives zero by multilinearity and alternation. Conversely extend an independent family to a basis. In that basis define $\varphi$ to be the determinant of the first $k$ coordinates of its arguments. It is alternating multilinear and takes value one on $v_1,\ldots,v_k$. Part (b) then makes their wedge nonzero. This uses no later Laplace expansion.
\end{proof}

The universal property also removes the original basis choice: two constructions with that property have mutually inverse maps fixed on all decomposable wedges.

\begin{definition}\label{linear:def:compound}
For linear $f:V\to W$, its \emph{exterior power} $\Lambda^kf$ is specified by
\[\Lambda^kf(v_1\wedge\cdots\wedge v_k)=f(v_1)\wedge\cdots\wedge f(v_k).
\]
The universal property supplies the unique linear map. For $A\in M_{m,n}(\mathbb K)$ its \emph{compound matrix} is $C_k(A)=(\det A_{I,J})_{|I|=|J|=k}$, with all index sets increasing.
\end{definition}

\begin{proposition}\label{linear:prop:compound}
A matrix $A$ for $f$ induces matrix $C_k(A)$ for $\Lambda^kf$ in the wedge bases. Exterior powers preserve composition and identity:
\[\Lambda^k(g\circ f)=\Lambda^kg\circ\Lambda^kf,\qquad\Lambda^k\mathrm{id}=\mathrm{id}.
\]
For an endomorphism of an $n$-dimensional space, $\Lambda^nf$ acts on its one-dimensional space by the scalar $\det f$.
\end{proposition}
\begin{proof}
The image of $e_J$ is the wedge of the $J$ columns of $A$, whose $I$ coefficient is $\det A_{I,J}$. Composition and identity agree on decomposable wedges, which span the exterior space. At top degree only one minor remains, the full determinant.
\end{proof}

Thus determinant multiplicativity is composition on the top exterior line. Keeping lower-degree coordinates gives the 1812 identity independently found by Binet and Cauchy.

\begin{theorem}[Cauchy--Binet]\label{linear:thm:cauchybinet}
For $A\in M_{m,n}(\mathbb K)$ and $B\in M_{n,p}(\mathbb K)$, one has $C_k(AB)=C_k(A)C_k(B)$, or entrywise
\[\det(AB)_{I,J}=\sum_{|S|=k}\det A_{I,S}\det B_{S,J}.
\]
In particular, for $A\in M_{k,n}$ and $B\in M_{n,k}$,
\[\det(AB)=\sum_{|S|=k}\det A_{[k],S}\det B_{S,[k]},\]
with an empty sum equal to zero when $n<k$.
\end{theorem}
\begin{proof}
Apply the composition proposition to $\mathbb K^p\xrightarrow B\mathbb K^n\xrightarrow A\mathbb K^m$. Taking its matrix gives the formula; if $k>n$, the intermediate exterior space is zero.
\end{proof}

\begin{corollary}[Pythagoras]\label{linear:cor:pythagoras}
For real $M\in M_{n,k}$,
\[\det(M^{\mathsf T}M)=\sum_{|I|=k}\det(M_I)^2.
\]
Thus squared $k$-volume equals the sum of squared projected $k$-volumes on the coordinate $k$-planes.
\end{corollary}
\begin{proof}
Use Cauchy--Binet on $M^{\mathsf T}M$. The algebraic identity is proved here; identifying the square root of the Gram determinant with geometric $k$-volume is established in Appendix~\ref{linear:ch:inner}. For $k=1$ this is ordinary Pythagoras. At $k=2,n=3$, division by four gives the corresponding triangle-area identity, related to de Gua's 1783 theorem.
\end{proof}

\section{Minors of the inverse and Dodgson's condensation}\label{linear:sec:condensation}

Write $A_{I,J}$ for selected rows and columns, $I^c$ for their complement in $[n]$, and $s(I)=\sum_{i\in I}i$. Jacobi's 1833 identity extends the inverse cofactor formula from single entries to all minors.

\begin{theorem}[Jacobi's complementary minors]\label{linear:thm:jacobi}
If $A\in GL_n(\mathbb K)$ and $|I|=|J|=k$, then
\[\det(A^{-1})_{J,I}=(-1)^{s(I)+s(J)}\frac{\det A_{I^c,J^c}}{\det A}.\]
\end{theorem}
\begin{proof}
First take $I=J=[k]$ and block the matrices as $A=\begin{pmatrix}P&Q\\R&S\end{pmatrix}$ and $A^{-1}=\begin{pmatrix}W&X\\Y&Z\end{pmatrix}$. Matrix multiplication gives
\[A\begin{pmatrix}W&0\\Y&I\end{pmatrix}=\begin{pmatrix}I&Q\\0&S\end{pmatrix}.
\]
Taking block-triangular determinants yields $\det A\det W=\det S$, without assuming $S$ invertible.

For general $I,J$, move their indices to the front by order-preserving row and column permutations. The inverse has its corresponding $J,I$ block at the front, so the preceding case applies. Moving $I=\{i_1<\cdots<i_k\}$ requires $\sum_r(i_r-r)=s(I)-k(k+1)/2$ adjacent swaps, and similarly for $J$. Their combined sign is $(-1)^{s(I)+s(J)-k(k+1)}=(-1)^{s(I)+s(J)}$, proving the identity. Empty and full index sets agree with the empty-determinant convention.
\end{proof}

Let $A^i_j$ delete row $i$ and column $j$, and let $A^{1,n}_{1,n}$ delete both corner rows and columns. The $k=1$ identity says $(A^{-1})_{ji}=(-1)^{i+j}\det A^i_j/\det A$. Define $\operatorname{adj}(A)$ by these transposed cofactors; its polynomial version survives without an inverse.
\begin{corollary}[adjugate identity]\label{linear:cor:adjugate}
For a square matrix over any commutative ring $R$,
\[A\operatorname{adj}(A)=\operatorname{adj}(A)A=(\det A)I.\]
\end{corollary}
\begin{proof}
Each entry of the difference is an integer-coefficient polynomial $p$ in the matrix entries. It vanishes on invertible rational matrices by Jacobi. Thus $p\det$ vanishes at every rational matrix, including singular ones. A polynomial over an infinite field vanishing everywhere is zero, by induction on its variables. Since the integer polynomial ring is an integral domain, $p=0$. Substitution proves the identity over any commutative ring, and the same argument covers multiplication on either side.
\end{proof}

Appendix~\ref{linear:ch:eigen} applies this over $\mathbb K[X]$ for Cayley--Hamilton. Two corner rows and columns give another useful identity, found by Desnanot in small sizes in 1819 and by Jacobi in general.

\begin{theorem}[Desnanot--Jacobi identity]\label{linear:thm:desnanot}
For $A\in M_n(\mathbb K)$ with $n\ge2$,
\[\det A\det A^{1,n}_{1,n}=\det A^1_1\det A^n_n-\det A^1_n\det A^n_1,\]
where the empty central minor has determinant one at $n=2$.
\end{theorem}
\begin{proof}
For invertible $A$, Jacobi with $I=J=\{1,n\}$ expresses the corner $2\times2$ inverse minor as $\det A^{1,n}_{1,n}/\det A$. Substituting the four cofactor entries of the inverse gives the numerator on the right divided by $(\det A)^2$. Multiply through to obtain the identity.

For arbitrary $A$ over any field, work over $\mathbb K(t)$. The matrix $A+tI$ is invertible there because its determinant is monic of degree $n$ in $t$. Apply the identity, now an equality of polynomials in $t$, and evaluate at zero. This argument also works when $\mathbb K$ is finite.
\end{proof}

\begin{corollary}[Dodgson's condensation]\label{linear:cor:dodgson}
For $0\le k\le n$, let $A^{(k)}$ contain the contiguous $k\times k$ minors of $A$, in an $(n-k+1)\times(n-k+1)$ array. Thus $A^{(0)}$ is all ones, $A^{(1)}=A$, and $A^{(n)}=(\det A)$. For $1\le k<n$ and $1\le i,j\le n-k$,
\[A^{(k+1)}_{ij}A^{(k-1)}_{i+1,j+1}
 =A^{(k)}_{ij}A^{(k)}_{i+1,j+1}-A^{(k)}_{i,j+1}A^{(k)}_{i+1,j}.\]
\end{corollary}
\begin{proof}
Apply Desnanot--Jacobi to the contiguous $(k+1)$ block at $(i,j)$. Its four corner minors are the right-hand entries, and its central minor is the second factor on the left.
\end{proof}

When all required divisors are nonzero, this recurrence computes successive arrays using only $2\times2$ determinants and division. Over integers these divisions are exact because their quotients are minors. A zero divisor can stop the algorithm even for an invertible matrix. Row permutations sometimes fix it. For a universal symbolic repair, replace each entry $a_{ij}$ by $a_{ij}+t_{ij}$, with independent indeterminates, and work in their rational-function field. Every required minor is then a nonzero polynomial, since its highest-degree part is a generic determinant. The final result simplifies to the determinant polynomial, where setting every $t_{ij}=0$ is valid. Perturbing only by $tI$ need not fix off-diagonal zero divisors; an added exercise exhibits this distinction.

\begin{remark}
Robbins--Rumsey's $\lambda$-determinant changes the subtraction to addition with coefficient $\lambda$. Its expansion is indexed by alternating sign matrices. The Mills--Robbins--Rumsey conjecture counts those matrices by $\prod_{j=0}^{n-1}(3j+1)!/(n+j)!$. Zeilberger's 1996 proof and Kuperberg's six-vertex-model proof establish this quoted enumeration. Bressoud's book explains the history and mathematics; neither proof is a premise of condensation here.
\end{remark}

\section{Paths that never meet}\label{linear:sec:lgv}

Let a finite directed acyclic graph have edge weights in a commutative ring. Give each directed path the product of its edge weights, and put
\[e(a,b)=\sum_{P:a\to b}w(P),\]
including the length-zero path of weight one when $a=b$. For sources $A_i$ and sinks $B_j$, a path system consists of a permutation $\sigma$ and paths $P_i:A_i\to B_{\sigma(i)}$. Its sign is $\varepsilon(\sigma)$ and its weight is the product of path weights. Call it \emph{non-intersecting} when distinct paths share no vertex.

\begin{theorem}[Lindstr\"om, 1973; Gessel--Viennot, 1985]\label{linear:thm:lgv}
The matrix $M=(e(A_i,B_j))$ satisfies
\[\det M=\sum_{\text{non-intersecting }(\sigma,P)}\varepsilon(\sigma)\prod_iw(P_i).\]
\end{theorem}
\begin{proof}
Expanding each product in the permutation formula sums signed weights over all path systems. We cancel the intersecting ones in pairs.

Choose a topological ordering of the finite acyclic graph: repeatedly remove a vertex with no incoming edge; such a vertex must exist, since following incoming edges forever in a finite graph would create a cycle. In an intersecting system, take the earliest shared vertex $v$ in this ordering, then the two smallest path indices passing through it. Swap their tails after $v$. Acyclicity ensures the new walks have no repeated vertices and are paths. The edge multiset and weight stay the same, while the two endpoints are exchanged, reversing the permutation sign.

This choice is stable under the swap. All changed vertices after $v$ are later in the topological order, and every prefix through $v$ is unchanged. Thus $v$ is still the earliest shared vertex and the same indices are chosen. Swapping twice restores the system. Each paired contribution sums to zero, also in characteristic two, leaving exactly the stated non-intersecting sum. Figure~\ref{linear:fig:lgv} illustrates the move.
\end{proof}
\begin{figure}[ht]
\centering
\fig{linear-ch02-lgv}
\caption{The tail swap. Two paths meet at the marked vertex; exchanging what comes after it changes the permutation by a
transposition and keeps the edges, so the two systems cancel in the determinant.}
\label{linear:fig:lgv}
\end{figure}

When only the identity pairing permits disjoint paths, the determinant becomes their unsigned weighted count. Ordered lattice endpoints provide a useful condition. North-east paths below use unit steps east and north; any fixed endpoint family lies in a finite bounding rectangle, so the finite-graph theorem applies.

\begin{lemma}\label{linear:lem:planar}
Suppose lattice sources $A_1,\ldots,A_k$ and sinks $B_1,\ldots,B_k$ each move weakly north-west with increasing index: $x$ does not increase and $y$ does not decrease, with consecutive points distinct. For $i<j$ and $l>m$, every north-east path $A_i\to B_l$ meets every north-east path $A_j\to B_m$. Thus a non-intersecting system must have identity permutation.
\end{lemma}
\begin{proof}
A north-east path visits exactly one vertex on each integer anti-diagonal $s=x+y$ between its endpoints. Write its horizontal coordinate as $x_P(s)$. Successive coordinates change by zero or one.

For paths $P:A_i\to B_l$ and $Q:A_j\to B_m$, their common anti-diagonals range from $s_0=\max(s(A_i),s(A_j))$ to $s_1=\min(s(B_l),s(B_m))$. This range is nonempty: the north-west endpoint inequalities and existence of both paths give $s(A_i)\le s(B_m)$ and $s(A_j)\le s(B_l)$, in addition to the direct endpoint inequalities.

At $s_0$, $x_P-x_Q\ge0$. If $Q$ starts there, use $x(A_j)\le x(A_i)$; if $P$ starts there, use $y_Q\ge y(A_j)\ge y(A_i)$. At $s_1$, $x_P-x_Q\le0$. If $Q$ ends there, use $x_P\le x(B_l)\le x(B_m)$; if $P$ ends there, use $y_Q\le y(B_m)\le y(B_l)$. This integer difference changes by at most one per anti-diagonal, so it vanishes somewhere. The two vertices there coincide. Finally every nonidentity permutation has an inversion, so the first assertion excludes it in a disjoint system.
\end{proof}

A path from $(x,y)$ to $(x',y')$ has $\binom{(x'-x)+(y'-y)}{x'-x}$ choices when both increments are non-negative, and none otherwise. Thus ordered path counts produce binomial determinants. Use $\binom nk=0$ when $k<0$ or $k>n$ for $n\ge0$.

\begin{example}\label{linear:ex:binomialdet}
One obtains $\det(\binom{i+j}{i})_{0\le i,j\le n}=1$ by choosing $A_i=(-i,0)$ and $B_j=(0,j)$. The path count has the displayed binomial coefficient, and the planar lemma permits only identity pairing. There is exactly one disjoint system. The zeroth path is the single vertex $(0,0)$. Inductively the $i$th path must stay in column $-i$ until height $i$, since any earlier east step lands on the preceding hook. It then runs east to $(0,i)$. These nested hooks are disjoint and exhaust the possibilities, proving the determinant equals one.
\end{example}

\subsection*{Lozenge tilings}

Let $H(a,b,c)$ be the triangular-lattice hexagon with positive integer side lengths $a,b,c,a,b,c$ and angles $120^\circ$. Draw the length-$c$ sides vertically. A lozenge consists of two unit triangles sharing an edge. Of its three orientations, two have vertical sides; chains of those tiles will encode the tiling in Figure~\ref{linear:fig:lozenge}.
\begin{figure}[ht]
\centering
\fig{linear-ch02-lozenge}
\caption{A lozenge tiling of $H(4,5,3)$, seen as a pile of cubes in a $4\times5\times3$ box. The lozenges with vertical
edges (shaded) form three chains, one from each unit edge of the left vertical side to the right vertical side; their
midlines are drawn. The chains never meet, and they determine the tiling.}
\label{linear:fig:lozenge}
\end{figure}

\begin{theorem}[counting lozenge tilings]\label{linear:thm:lozenge}
The number of tilings of $H(a,b,c)$ is
\[T(a,b,c)=\det\left(\binom{a+b}{a+i-j}\right)_{1\le i,j\le c}.\]
\end{theorem}
\begin{proof}
Start at a unit edge on the left vertical side. The triangle inside belongs to a tile with two vertical sides: that boundary edge cannot be the edge joining its triangles. Leave through the other vertical side and continue. Each tile moves the chain strictly right, ending on the right vertical boundary. There are $c$ chains. Distinct chains cannot share a tile, because its two vertical sides determine its entry and exit. Following any such tile leftwards shows it belongs to one of these chains.

The chains also determine every remaining tile. A triangle outside them must be paired with the triangle across its vertical edge. That neighbour exists, since all vertical boundary edges belong to chains, and it is also outside the chains: a chain using the neighbour would continue across this edge into the first triangle. Thus the unused triangles pair uniquely.

Record a chain at the midpoints of its vertical sides. Its two step directions are $v$, down and right at $30^\circ$, and $u$, up and right at $30^\circ$. In coordinates relative to $v,u$, these are east and north, while the upward vertical displacement is $d=u-v=(-1,1)$. After translating the origin, the left midpoints are $A_i=(-i,i)$ and the right midpoints are $B_j=(a-j,b+j)$. Thus a path $A_i\to B_j$ has $a+i-j$ east steps out of $a+b$, giving the stated matrix entry.

Disjoint chains give vertex-disjoint lattice paths. Conversely, identity-paired north-east paths give disjoint chains of vertical-sided lozenges. They stay in the hexagon: each path is its starting midpoint plus at most $a$ steps in direction $v$ and at most $b$ in direction $u$, within the hexagon's three-segment description by the $v,u,d$ directions. The half-edge offsets at the vertical boundaries place the corresponding lozenges inside the same hexagon. Pairing every remaining triangle across its vertical edge then completes the tiling, as above. This construction reverses the chain extraction.

Both endpoint lists move north-west. The planar lemma permits only identity pairing, and LGV counts exactly these disjoint systems by the determinant.
\end{proof}

At $a=b=c=2$ the determinant is $\binom42^2-\binom43\binom41=36-16=20$. There is also a product evaluation, older than the path lemma.

\begin{remark}[MacMahon's formula]
MacMahon's boxed-plane-partition enumeration, developed in 1896--1916, gives
\[T(a,b,c)=\prod_{i=1}^a\prod_{j=1}^b\prod_{k=1}^c\frac{i+j+k-1}{i+j+k-2}.
\]
We quote the general evaluation; Krattenthaler's survey proves the corresponding binomial determinant formula. The cases $c=1,2$ are proved in Exercise~\ref{linear:ex:hex2}, not inferred from numerical checks. For $a=b=c=3$ the product is $980$. The cube-stack interpretation in Figure~\ref{linear:fig:lozenge} identifies these tilings with plane partitions in a box. Cohn, Larsen and Propp's 1998 limit-shape theorem describes frozen corners and a liquid region whose boundary is a circle for a large regular hexagon. This probabilistic theorem is further context, not part of the determinant proof.
\end{remark}

\begin{beauty}
For path-count matrices, the permutation expansion has a concrete cancellation. One tail swap pairs every intersecting configuration with a contribution of opposite sign. The survivors are exactly the disjoint configurations. Karlin--McGregor's 1959 probability formulas use a related mechanism for noncolliding walks, one connection between determinants and the random-matrix theory of Volume~12.
\end{beauty}

\begin{figure}[ht]
\centering\fig{appB-rw-minors}
\caption{For the columns of $M$, the three projected signed areas are $1,1,-1$. Cauchy--Binet adds their squares to the Gram determinant $3$, so the parallelogram has area $\sqrt3$.}
\label{fig:v8appBminors}
\end{figure}
\begin{figure}[ht]
\centering\fig{appB-rw-condense}
\caption{Condensation multiplies the larger minor by the central smaller minor before equating it to a difference of products. The identity holds at a zero central minor, but division there does not determine the larger minor.}
\label{fig:v8appBcondense}
\end{figure}
\section{Exercises}
\begin{exercise}[$\star$]\label{linear:ex:latticetriangle}
Use shoelace to find the area with consecutive vertices $(0,0),(4,0),(5,3),(2,5),(-1,2)$. Prove that every integer-coordinate triangle has half-integer area.
\end{exercise}
\begin{exercise}[$\star$]\label{linear:ex:laplace}
For $A\in M_n(\mathbb K)$ and $|J|=k$, prove the column-set Laplace expansion
\[\det A=\sum_{|I|=k}(-1)^{s(I)+s(J)}\det A_{I,J}\det A_{I^c,J^c}.\]
Factor the alternating dependence on the columns in $J$ through an exterior power.
\end{exercise}
\begin{exercise}[$\star\star$]\label{linear:ex:lagrangeid}
(a) Derive $\|a\|^2\|b\|^2-(a\cdot b)^2=\sum_{i<j}(a_ib_j-a_jb_i)^2$ from Cauchy--Binet. Deduce Cauchy--Schwarz and its exact equality condition, including zero vectors. (b) For $v,w\in\R^3$, obtain $\|v\times w\|^2=\|v\|^2\|w\|^2-(v\cdot w)^2$.
\end{exercise}
\begin{exercise}[$\star\star$]\label{linear:ex:charpoly}
Expand $\det(XI-A)$ to prove
\[\det(XI-A)=\sum_{k=0}^n(-1)^k\operatorname{tr}(C_k(A))X^{n-k}.\]
Thus each coefficient is a signed sum of principal minors.
\end{exercise}
\begin{exercise}[$\star$]\label{linear:ex:compoundeig}
For upper triangular $A$, prove that $C_k(A)$ is upper triangular in lexicographic subset order, with diagonal $\prod_{i\in I}a_{ii}$. Using the earlier complex triangularisation in Chapter~\ref{linear:ch:spectral}, deduce that $C_k(A)$ has eigenvalues $\lambda_{i_1}\cdots\lambda_{i_k}$ for all increasing index choices, with multiplicities.
\end{exercise}
\begin{exercise}[$\star\star$]\label{linear:ex:adjadj}
Prove $\operatorname{adj}(\operatorname{adj}A)=(\det A)^{n-2}A$ for every $n\ge2$, including singular matrices.
\end{exercise}
\begin{exercise}[$\star$]\label{linear:ex:dodgson}
Condense $\begin{pmatrix}1&2&3&4\\2&1&2&3\\3&2&1&2\\4&3&2&1\end{pmatrix}$ exactly by hand. Find an invertible real $3\times3$ matrix where a zero centre stops condensation and give a valid repair.
\end{exercise}
\begin{exercise}[$\star\star$]\label{linear:ex:plucker}
For any $4\times2$ matrix, let $p_{ij}$ denote its row-pair minors. Prove $p_{12}p_{34}-p_{13}p_{24}+p_{14}p_{23}=0$ over any field. Conversely realise every nonzero six-tuple satisfying it as the minors of a rank-two matrix, retaining the tuple's original scale. Over $\R$, this represents two-dimensional subspaces of $\R^4$ by points of a quadric in projective five-space; equivalently these are projective lines in projective three-space.
\end{exercise}
\begin{exercise}[$\star\star$]\label{linear:ex:catalan}
The Catalan number $C_m=(m+1)^{-1}\binom{2m}m$ counts north-east paths $(0,0)\to(m,m)$ staying below the diagonal. Prove both $\det(C_{i+j})_{0\le i,j\le n}=1$ and $\det(C_{i+j+1})_{0\le i,j\le n}=1$ for every $n\ge0$.
\end{exercise}
\begin{exercise}[$\star\star$]\label{linear:ex:hex2}
Find a direct count $T(a,b,1)=\binom{a+b}a$ and a closed expression for $T(a,b,2)$. Prove agreement with MacMahon's product in these two cases.
\end{exercise}
\begin{exercise}[$\star\star$]\label{linear:ex:vandermondelgv}
For increasing non-negative integers $a_1<\cdots<a_k$ and $b_1<\cdots<b_k$, prove $\det(\binom{a_i}{b_j})\ge0$, and show it vanishes exactly when some $a_i<b_i$.
\end{exercise}
\begin{exercise}[$\star\star$]\label{linear:ex:calisson}
Prove that every tiling of $H(a,b,c)$ contains $ab,bc,ca$ lozenges in its three orientations. Thus a regular side-$n$ hexagon always has $n^2$ of each kind, independently of its tiling.
\end{exercise}

\begin{exercise}[$\star$]\label{ex:v8cBa}
\emph{A projected-area calculation.} For $M=\begin{pmatrix}1&0\\0&1\\1&1\end{pmatrix}$, compute every $2\times2$ minor and $\det(M^{\mathsf T}M)$. Find the parallelogram area over $\R$. Over $\mathbb F_3$, show that its columns remain independent although the same Gram determinant is zero. Why does that not contradict the real volume identity?
\end{exercise}
\begin{exercise}[$\star\star$]\label{ex:v8cBb}
\emph{A diagonal perturbation can still fail.} Run the formal condensation setup on $I_4+tI_4$. Identify a required division with zero denominator. Explain why replacing each entry of $I_4$ by $a_{ij}+t_{ij}$, using independent indeterminates, removes every such obstruction over the rational-function field.
\end{exercise}
\begin{exercise}[$\star\star$]\label{ex:v8cBc}
\emph{A binomial Vandermonde determinant.} For distinct real $a_1,\ldots,a_k$, interpret $\binom{x}{j}=x(x-1)\cdots(x-j+1)/j!$ as a polynomial and prove
\[\det\left(\binom{a_i}{j-1}\right)_{1\le i,j\le k}
 =\frac{\prod_{i<l}(a_l-a_i)}{\prod_{j=0}^{k-1}j!}.\]
Use this to prove strict positivity for increasing non-negative integer $a_i$ in this special case of Problem~\ref{linear:ex:vandermondelgv}.
\end{exercise}
\begin{exercise}[$\star\star$]\label{ex:v8cBd}
\emph{A tiling recurrence from condensation.} Set $T(a,b,0)=1$. For positive $a,b$ and $c\ge2$, apply Desnanot--Jacobi to the tiling determinant to prove
\[T(a,b,c)T(a,b,c-2)=T(a,b,c-1)^2
 -T(a+1,b-1,c-1)T(a-1,b+1,c-1).\]
Here a zero side is allowed in the two latter terms and gives one. Justify this boundary convention from the determinant, and verify the recurrence at $a=b=c=2$ using exact small determinants.
\end{exercise}

\section{Hints}
\begin{hint}{linear:ex:latticetriangle}
The five determinants sum to an even integer. For a general lattice triangle, translate one vertex to zero.
\end{hint}
\begin{hint}{linear:ex:laplace}
Freeze the other columns and evaluate the induced form on each coordinate wedge. Bringing the chosen rows and columns forward determines the sign.
\end{hint}
\begin{hint}{linear:ex:lagrangeid}
Use the matrix with columns $a,b$ and identify its Gram determinant.
\end{hint}
\begin{hint}{linear:ex:charpoly}
In each column choose either $Xe_j$ or $-A_j$, then group choices by subset size.
\end{hint}
\begin{hint}{linear:ex:compoundeig}
Every basis vector occurring in $Ae_i$ has index at most $i$. Exterior functoriality preserves similarity.
\end{hint}
\begin{hint}{linear:ex:adjadj}
Prove it first using inverses, then replace $A$ by $A+tI$ over $\mathbb K(t)$ for the polynomial-identity extension.
\end{hint}
\begin{hint}{linear:ex:dodgson}
The first minor array has diagonal $-3$ and off-diagonal entries one. A row swap repairs a suitable $3\times3$ example; independent generic entry perturbations give a universal symbolic repair.
\end{hint}
\begin{hint}{linear:ex:plucker}
Prove the relation by direct row-coordinate cancellation, without dividing by two. When $p_{12}=d\ne0$, start with rows $(1,0),(0,d)$.
\end{hint}
\begin{hint}{linear:ex:catalan}
Use the region $y\le x$ with sources $(-i,-i)$ and sinks $(j,j)$ or $(j+1,j+1)$. Show disjoint paths are forced to make nested corners.
\end{hint}
\begin{hint}{linear:ex:hex2}
Expand the size-two determinant and telescope the product separately at heights one and two.
\end{hint}
\begin{hint}{linear:ex:vandermondelgv}
Take paths from $(0,-a_i)$ to $(b_j,-b_j)$. If all $a_i\ge b_i$, construct disjoint east-then-north corners.
\end{hint}
\begin{hint}{linear:ex:calisson}
Each chain contains exactly $a$ tiles of one kind and $b$ of another. Count total unit triangles for the remainder.
\end{hint}
\begin{hint}{ex:v8cBa}
The projected minors are $1,1,-1$. In characteristic three, a nonzero sum of squares may vanish.
\end{hint}
\begin{hint}{ex:v8cBb}
At the third-minor stage, look at the denominator in row two, column three of the original matrix.
\end{hint}
\begin{hint}{ex:v8cBc}
Changing from monomials to binomial polynomials is triangular with diagonal $1/j!$.
\end{hint}
\begin{hint}{ex:v8cBd}
The diagonal corner minors keep $a$, while deleting opposite row and column corners changes $a$ to $a+1$ or $a-1$ and keeps $a+b$ fixed.
\end{hint}

\section{Solutions}
\begin{solution}{linear:ex:latticetriangle}
\idea Integer coordinates make twice the oriented area an integer determinant sum.
\writeup The five consecutive determinants are $0,12,19,9,0$, with total $40$. The positive orientation gives area $20$. For any lattice triangle, translation gives area $\tfrac12|\det(q-p,r-p)|$, a half-integer. Pick's theorem in Volume~9, Chapter~9, adds a boundary-and-interior interpretation but is unnecessary here.
\end{solution}
\begin{solution}{linear:ex:laplace}
\idea Freeze the unselected columns, then evaluate the remaining alternating form on coordinate wedges.
\writeup Let $\varphi(x_1,\ldots,x_k)$ be the determinant with these vectors inserted in columns $J$, the others fixed. Exterior factorisation gives $\det A=\sum_I\det A_{I,J}\tilde\varphi(e_I)$. To compute the latter coefficient, move rows $I$ and columns $J$ to the front, retaining their internal order. Their total swap parity is $s(I)+s(J)-k(k+1)$. The resulting matrix has block form $\begin{pmatrix}I_k&*\\0&A_{I^c,J^c}\end{pmatrix}$, proving $\tilde\varphi(e_I)=(-1)^{s(I)+s(J)}\det A_{I^c,J^c}$. Substitution proves the formula.
\end{solution}
\begin{solution}{linear:ex:lagrangeid}
\idea Two columns turn the Gram determinant into a sum of squared two-row minors.
\writeup (a) For $M=(a\mid b)$, its Gram matrix is $\begin{pmatrix}\|a\|^2&a\cdot b\\a\cdot b&\|b\|^2\end{pmatrix}$. Cauchy--Binet gives the claimed identity. Its non-negative right side proves Cauchy--Schwarz. Equality means every minor vanishes, equivalently rank at most one. Thus the vectors are linearly dependent, including when either is zero; if one is nonzero, the other is its scalar multiple. (b) In three dimensions the minors are exactly the cross-product coordinates up to signs, giving the second identity.
\end{solution}
\begin{solution}{linear:ex:charpoly}
\idea Choose one of the two summands in each column and collect the resulting principal minors.
\writeup Expanding columns $Xe_j-A_j$ gives a term for each subset $S$ of columns chosen from $-A$. The remaining identity columns force the unused row indices, leaving $(-1)^{|S|}X^{n-|S|}\det A_{S,S}$. Summing over $|S|=k$ gives $(-1)^kX^{n-k}\operatorname{tr}(C_k(A))$, including $k=0$ under the empty-minor convention.
\end{solution}
\begin{solution}{linear:ex:compoundeig}
\idea A triangular map can only lower wedge indices; its diagonal selections multiply.
\writeup Write $Ae_i=a_{ii}e_i+w_i$ with $w_i\in\operatorname{span}(e_1,\ldots,e_{i-1})$. Expanding a wedge gives its original $e_I$ with coefficient $\prod_{i\in I}a_{ii}$. Every other nonzero term has an increasing index list componentwise at most $I$, and strictly smaller somewhere, hence lexicographically earlier. This proves triangularity and its diagonal. For complex $A=PTP^{-1}$, the earlier triangularisation and exterior composition give $C_k(A)=C_k(P)C_k(T)C_k(P)^{-1}$. Its eigenvalues are therefore the diagonal products from $T$, counted by the index subsets. This argument uses the already proved main-chapter triangularisation, rather than a later appendix.
\end{solution}
\begin{solution}{linear:ex:adjadj}
\idea Prove a polynomial identity on invertible matrices and then specialise a formal invertible perturbation.
\writeup For invertible $A$, $\operatorname{adj}A=(\det A)A^{-1}$ and $\det(\operatorname{adj}A)=(\det A)^{n-1}$. Applying the same adjugate formula again gives $(\det A)^{n-2}A$. Both sides are polynomial because $n\ge2$. Replace an arbitrary $A$ by $A+tI$ over $\mathbb K(t)$, where it is invertible; the resulting polynomial identity at $t=0$ proves the claim, even over finite fields.
\end{solution}
\begin{solution}{linear:ex:dodgson}
\idea Keep each contiguous-minor array exact and track which interior entry divides the next stage.
\writeup Direct two-by-two determinants give
\[A^{(2)}=\begin{pmatrix}-3&1&1\\1&-3&1\\1&1&-3\end{pmatrix},\qquad
A^{(3)}=\begin{pmatrix}8&2\\2&8\end{pmatrix}.
\]
For example the top-left third minor is $(9-1)/1=8$ and the adjacent one is $(1+3)/2=2$. The final determinant is $(64-4)/(-3)=-20$.

The real matrix $B=\begin{pmatrix}1&1&0\\1&0&1\\0&1&1\end{pmatrix}$ has determinant $-2$ but central entry zero, so the last division fails. Swap its first two rows: the new centre is one, condensation succeeds, and one sign change recovers the original determinant. More generally the independent-entry perturbation described after Corollary~\ref{linear:cor:dodgson} works over a rational-function field. A diagonal perturbation alone is insufficient in all sizes, as Exercise~\ref{ex:v8cBb} shows.
\end{solution}
\begin{solution}{linear:ex:plucker}
\idea A two-coordinate row identity proves the relation without characteristic restrictions; a nonzero minor then provides coordinates for reconstruction.
\writeup Let $r_i$ be the two-component rows. Expanding their two coordinates gives $p_{23}r_1-p_{13}r_2+p_{12}r_3=0$. Take its determinant with $r_4$. Bilinearity gives $p_{23}p_{14}-p_{13}p_{24}+p_{12}p_{34}=0$, with no division by two.

Conversely assume $d=p_{12}\ne0$. Use rows
\[(1,0),\quad(0,d),\quad(-p_{23}/d,p_{13}),\quad(-p_{24}/d,p_{14}).\]
The first five minors have exactly their prescribed values. The last is $(p_{13}p_{24}-p_{23}p_{14})/d=p_{34}$ by the relation. If another minor is nonzero, permute rows to put that pair first, keeping the induced signs on every prescribed coordinate, perform the construction and undo the permutation. A nonzero minor ensures rank two. Scaling all six coordinates represents the same two-plane, while the construction itself has retained the original scale.
\end{solution}
\begin{solution}{linear:ex:catalan}
\idea Catalan entries count paths constrained below the diagonal, and disjoint families are forced into nested corners.
\writeup Work in a finite bounding region of $y\le x$, with $A_i=(-i,-i)$ and $B_j=(j,j)$. Translation gives path count $C_{i+j}$. Any inverted pairing $i<j$, $l>m$ forces intersection: at anti-diagonal $s=-2i$, the path from $A_i$ has $x=-i$ and the other has $x\ge-i$; at $s=2m$ the path ending at $B_m$ has $x=m$ and the other has $x\ge m$. The difference changes by at most one, so it hits zero. Thus only identity pairing survives LGV.

The zeroth path is the single origin. Inductively, once earlier paths are the corners east from $(-k,-k)$ to $(k,-k)$ then north to $(k,k)$, path $i$ cannot turn north before $x=i$: it would enter the preceding corner. Since its endpoint has $x=i$, it must turn there, determining it uniquely. These corners are disjoint, so the first determinant is one.

For the shifted entries use $B_j=(j+1,j+1)$. The same anti-diagonal comparison excludes inverted endpoints. The base path is forced east then north by the diagonal constraint; induction forces the corners east from $(-i,-i)$ to $(i+1,-i)$ then north to $(i+1,i+1)$. There is again one system, giving the second determinant.
\end{solution}
\begin{solution}{linear:ex:hex2}
\idea Evaluate the first two determinant sizes and telescope the matching product factors.
\writeup A single chain gives $T(a,b,1)=\binom{a+b}a$. The height-one product telescopes over $j$, then $i$, to $\prod_{i=1}^a(i+b)/i$, the same binomial coefficient. At height two,
\[T(a,b,2)=\binom{a+b}a^2-\binom{a+b}{a+1}\binom{a+b}{a-1}
 =\binom{a+b}a^2\frac{a+b+1}{(a+1)(b+1)}.
\]
Here the neighbouring binomial ratios are $b/(a+1)$ and $a/(b+1)$. The product's second-layer factors telescope to $\prod_{i=1}^a(i+b+1)/(i+1)=\binom{a+b}a(a+b+1)/((a+1)(b+1))$. Multiply by the first layer to prove agreement, including the value $20$ at $a=b=2$.
\end{solution}
\begin{solution}{linear:ex:vandermondelgv}
\idea Interpret the binomial entries as paths and construct a disjoint family whenever every diagonal entry is feasible.
\writeup Take $A_i=(0,-a_i)$ and $B_j=(b_j,-b_j)$. There are $b_j$ east steps and $a_i-b_j$ north steps, giving $\binom{a_i}{b_j}$. Reversing both endpoint lists makes each move north-west and changes the determinant by the square of a permutation sign, hence leaves it unchanged. The planar lemma and LGV make the determinant a non-negative count with identity pairing. If $a_i<b_i$, that diagonal path is impossible, so the count is zero.

Otherwise let each path go east to $(b_i,-a_i)$, then north to $B_i$. Distinct horizontal parts use different rows and vertical parts use different columns. For $i<j$, the $j$th horizontal segment lies below the $i$th vertical segment, while the $i$th horizontal segment stops before reaching the $j$th column. Therefore the paths are disjoint and contribute at least one, proving strict positivity.
\end{solution}
\begin{solution}{linear:ex:calisson}
\idea The two chain-step counts fix two orientations; the hexagon's area fixes the third.
\writeup Each of the $c$ chains has $a$ east and $b$ north steps, so it contains $a$ and $b$ lozenges of the two vertical-sided kinds. Their totals are $ca$ and $cb$. The hexagon splits into three lattice parallelograms, containing $2ab,2bc,2ca$ unit triangles. Every tiling therefore has $ab+bc+ca$ lozenges, leaving $ab$ of the third orientation. At equal side lengths all three totals are $n^2$. David--Tomei's cube-stack picture gives another view of these fixed projected face counts.
\end{solution}
\begin{solution}{ex:v8cBa}
\idea Over the reals a sum of squared minors measures area; in a finite field it need not be positive.
\writeup The minors are $p_{12}=1,p_{13}=1,p_{23}=-1$. The Gram matrix is $\begin{pmatrix}2&1\\1&2\end{pmatrix}$, with determinant three, equal to $1^2+1^2+(-1)^2$. Thus the real area is $\sqrt3$. Modulo three the columns remain independent because $p_{12}=1$, but the Gram determinant vanishes. Cauchy--Binet still holds as an algebraic identity; the real inference from a zero sum of squares to vanishing of every summand does not apply in $\mathbb F_3$.
\end{solution}
\begin{solution}{ex:v8cBb}
\idea Formal invertibility of the whole matrix does not guarantee that each interior divisor is nonzero.
\writeup The matrix $(1+t)I_4$ has zero entry in position $(2,3)$. Computing its contiguous third minor at $(1,2)$ requires division by exactly that entry, so the literal algorithm encounters $0/0$. With independent $t_{ij}$, every contiguous minor has a nonzero generic determinant as its highest-degree part. Each required divisor is consequently nonzero in the rational-function field, and condensation succeeds there. The result is the determinant polynomial before specialising the parameters to zero.
\end{solution}
\begin{solution}{ex:v8cBc}
\idea Binomial polynomials have the same degrees as monomials, with known leading coefficients.
\writeup The change from $1,x,\ldots,x^{k-1}$ to $\binom{x}{0},\ldots,\binom{x}{k-1}$ is triangular with diagonal $1/j!$. Evaluation therefore multiplies the ordinary Vandermonde determinant by $\prod_{j=0}^{k-1}(j!)^{-1}$, yielding the formula. Every factor $a_l-a_i$ is positive for increasing real nodes. For non-negative integer nodes this proves the desired special-case positivity directly; also $a_i\ge i-1=b_i$ in this case.
\end{solution}
\begin{solution}{ex:v8cBd}
\idea Removing adjacent corner rows and columns keeps or shifts the binomial parameter by one.
\writeup Let $M_{ij}=\binom{a+b}{a+i-j}$. Its two diagonal corner minors are both $T(a,b,c-1)$ and its central minor is $T(a,b,c-2)$. The opposite corner minors have lower parameters $a+1+i-j$ and $a-1+i-j$, giving $T(a+1,b-1,c-1)$ and $T(a-1,b+1,c-1)$. Desnanot--Jacobi proves the recurrence. When $a=0$ or $b=0$, the defining matrix is triangular with diagonal one, so its determinant is one; the size-zero case is the empty determinant. At $(2,2,2)$ the recurrence reads $20\cdot1=6^2-4\cdot4$, verified by the one- and two-dimensional determinants already evaluated.
\end{solution}

\section*{Sources and further reading}

\begin{itemize}
\item C.~L.~Dodgson, \emph{Condensation of determinants}, Proceedings of the Royal Society of London 15 (1866), 150--155;
and \emph{An Elementary Treatise on Determinants}, Macmillan, 1867.
\item D.~P.~Robbins and H.~Rumsey, \emph{Determinants and alternating sign matrices}, Advances in Mathematics 62 (1986),
169--184; W.~H.~Mills, D.~P.~Robbins and H.~Rumsey, \emph{Alternating sign matrices and descending plane partitions},
Journal of Combinatorial Theory A 34 (1983), 340--359.
\item D.~M.~Bressoud, \emph{Proofs and Confirmations: The Story of the Alternating Sign Matrix Conjecture}, Cambridge
University Press, 1999. A wonderful account of the whole story, with condensation at its start.
\item D.~Zeilberger, \emph{Proof of the alternating sign matrix conjecture}, Electronic Journal of Combinatorics 3(2)
(1996), R13; G.~Kuperberg, \emph{Another proof of the alternating-sign matrix conjecture}, International Mathematics
Research Notices 1996, 139--150.
\item B.~Lindstr\"om, \emph{On the vector representations of induced matroids}, Bulletin of the London Mathematical
Society 5 (1973), 85--90.
\item I.~Gessel and G.~Viennot, \emph{Binomial determinants, paths, and hook length formulae}, Advances in Mathematics 58
(1985), 300--321.
\item S.~Karlin and J.~McGregor, \emph{Coincidence probabilities}, Pacific Journal of Mathematics 9 (1959), 1141--1164.
\item M.~Aigner and G.~M.~Ziegler, \emph{Proofs from THE BOOK}, 6th ed., Springer, 2018, chapter ``Lattice paths and
determinants''.
\item P.~A.~MacMahon, \emph{Combinatory Analysis}, vol.~II, Cambridge University Press, 1916.
\item C.~Krattenthaler, \emph{Advanced determinant calculus}, S\'eminaire Lotharingien de Combinatoire 42 (1999), B42q.
Section~2 evaluates the binomial determinants behind MacMahon's formula.
\item H.~Cohn, M.~Larsen and J.~Propp, \emph{The shape of a typical boxed plane partition}, New York Journal of Mathematics
4 (1998), 137--165. The arctic circle for lozenge tilings.
\item G.~David and C.~Tomei, \emph{The problem of the calissons}, American Mathematical Monthly 96 (1989), 429--431: in
every lozenge tiling of a regular hexagon, the three kinds of lozenges appear equally often.
\item R.~Bhatia, \emph{Matrix Analysis}, Springer, 1997, Chapter~I, for exterior powers and compound matrices.
\end{itemize}
